下面希望利用 Banach 压缩原理, 证明 Banach 空间之间函数的隐函数定理. 为此, 首先需要定义导数的概念.

\begin{Definition}
\label{def:frechet-derivative}
设 $X$ 和 $Y$ 为 Banach 空间, $x_0$ 为 $X$ 中一点. 设 $F$ 为从 $x_0$ 的某个邻域映入 $Y$ 的映射. 若存在线性算子 $A\in\mathcal{L}(X,Y)$, 使得
\begin{equation}\label{eq:frechet-differentiability-linear-remainder-expansion}
F(x)=F(x_0)+A(x-x_0)+G(x),
\end{equation}
其中
\begin{equation}\label{eq:frechet-differentiability-vanishing-relative-remainder}
\lim_{x\to x_0}\|G(x)\|_Y/\|x-x_0\|_X=0.
\end{equation}
则称 $F$ 在 $x_0$ 处\textbf{可微}. 若这样的 $A$ 存在, 则称它为 $F$ 在 $x_0$ 处的 (Fréchet) 导数.
\end{Definition}

自然地, 若函数在定义域的所有点都可微, 就称该函数可微. 使用微积分中熟悉的记号, 如 $F^{\prime}(x_0)$ 或 $DF(x_0)$, 表示 Fréchet 导数. 若 $F^{\prime}(x)$ (作为 $\mathcal{L}(X,Y)$ 中的元素) 连续依赖于 $x$, 则称函数连续可微. 显然可以递归使用该定义来定义高阶导数. 例如, 若 $F^{\prime\prime}(x)$ 存在, 则它是 $\mathcal{L}(X,\mathcal{L}(X,Y))$ 中的元素. 也可以将 $F^{\prime\prime}(x)$ 看作从 $X$ 到 $Y$ 的双线性映射, 见 Problem~\ref{ex:continuous-bilinear-linear-operator-correspondence}. 导数的许多基本性质都可直接推广到 Fréchet 导数; 例如见 Problem~\ref{ex:frechet-derivative-chain-rule}.

\begin{Example}
\label{ex:pointwise-composition-frechet-derivative-example}
设 $f$ 为从 $\mathbb{R}$ 到 $\mathbb{R}$ 的连续可微函数, $X=Y=C([0,1])$. 映射 $f$ 通过逐点作用 $\phi(t)\mapsto f(\phi(t))$ 诱导从 $X$ 到 $Y$ 的映射. 作为从 $X$ 到 $Y$ 的映射, $f$ 可微, 且其在 $\phi$ 处的导数为将函数 $\psi(t)$ 映到 $f^{\prime}(\phi(t))\psi(t)$ 的线性映射. 见 Problem~\ref{ex:pointwise-composition-frechet-derivative}.
\end{Example}
